\section{Related Work}\label{sec:related_work}

Commercial systems for sewer and drainage networks provide resources for modeling, simulation, documentation, and operation.
Bentley OpenFlows Sewer, for instance, is a professional suite for modeling and managing sanitary sewer systems, with geospatial data integration and hydraulic simulation tools \cite{bentley2024}.
Such systems are robust, but they are usually directed toward complete professional environments and more detailed stages of the project lifecycle.

Urbano Hydra provides productivity resources for network projects, with synchronization between plans, profiles, and data tables \cite{peterschinegg}.
This integration between spatial representation and tabular data is relevant to \method because it reinforces the importance of keeping the network drawing, its attributes, and processing results consistent.
However, \method prioritizes the generation and preliminary analysis of routes from public data and topography instead of starting from an already consolidated technical drawing.

In the Brazilian context, GeoSan is a reference for the management and technical registry of water and sewer networks \cite{nexusgeosan}.
Its relationship with \method is complementary: while registry systems focus on the operation and record of existing networks, \method acts earlier, in the preliminary conception and feasibility analysis of new expansions.
The platform also differs by emphasizing interactive graph editing before processing, keeping the specialist inside the computational workflow.

Academic works that apply genetic algorithms, optimization models, and search techniques to sewer network layout are also relevant \cite{haghighi2012,weng2014,rodrigues2020}.
These studies show that a computational formulation of the problem can reduce manual effort and support the search for alternatives.
\method is positioned in this context as an integration-oriented solution: besides the routing algorithm, it incorporates geospatial data, elevation, persistence, interactive editing, and result visualization in a web application.
